On place dans un ballon $\qty{5.3}{\g}$ de carbonate de sodium
solide \ch{Na2CO3} auquel on ajoute un volume $V=\qty{50}{\mL}$ d'une solution
d'acide chlorhydrique de concentration $c=\qty{0.8}{\mol\per\L}$. Il se forme du
dioxyde de carbone, de l'eau et des ions sodium en solution.
L'équation de la réaction est~:
\[
    \ch{Na2CO3_{(s)} + 2 H^+_{(aq)} -> CO2_{(g)} + 2 Na^+_{(aq)} + H2O_{($\ell$)}}
\]

\begin{enumerate}
    \item Calculer les quantités de matière initiales des
          réactifs.~\textbf{(1pt)}
    \item Dresser le tableau d'avancement de la réaction. L'eau est en
          excès.~\textbf{(1pt)}
    \item Calculer l'avancement maximal $x_{\max}$ et déduire le réactif
          limitant.~\textbf{(1pt)}
    \item En déduire, dans l'état final~:
          \begin{enumerate}
              \item la composition en quantité de matière du
                    système~;~\textbf{(2pt)}
              \item le volume de gaz obtenu~;~\textbf{(1pt)}
              \item la concentration molaire des ions
                    \ch{Na^+_{(aq)}}~;~\textbf{(1pt)}
              %\item la masse de carbonate de sodium \ch{Na2CO3}.~\textbf{(1pt)}
          \end{enumerate}
\end{enumerate}
On supposera que, pendant la réaction, le volume de la solution ne varie pas et
que la pression égale à l'intérieur et à l'extérieur du ballon.
\begin{donnees}
    \begin{itemize}
        \item Masses molaires~: $M(\ch{Na})=\qty{23.0}{\g\per\mol}$~;
              $M(\ch{C})=\qty{12.0}{\g\per\mol}$~;
              $M(\ch{O})=\qty{16.0}{\g\per\mol}$.
        \item Volume molaire~: $V_{m}=\qty{24}{\L\per\mol}$.
    \end{itemize}
\end{donnees}

